\documentclass{article}
\usepackage[T1]{fontenc}
\usepackage{lmodern}
\usepackage{graphicx}
\usepackage{amsmath,amssymb}
\usepackage[a4paper,margin=2cm]{geometry}
\usepackage[absolute,overlay]{textpos}
\setlength{\TPHorizModule}{1mm}
\setlength{\TPVertModule}{1mm}
\usepackage[indonesian]{babel}
\usepackage{hyperref}
\setlength{\emergencystretch}{2em}
% unit: Halaman 5

\begin{document}

% === Halaman 5 (llm) ===

\textbf{Contoh:} Misalkan plainteks adalah \textcolor{red}{1101110100010001000101000101}, dibagi menjadi tiga buah blok masing-masing berukuran 8 bit:

\begin{center}
\textcolor{red}{11011101} \quad \textcolor{red}{00010001} \quad \textcolor{red}{01000101} \\
P1 \qquad\qquad P2 \qquad\qquad P3
\end{center}

\vspace{5mm}

Kunci yang digunakan adalah $\text{K} = \textcolor{blue}{01010110}$

Enkripsi blok pertama, \textcolor{red}{11011101}, dengan fungsi $\text{E}$ sebagai berikut:
\begin{enumerate}
  \item $\text{P1} \oplus \text{K} = \textcolor{red}{11011101} \oplus \textcolor{blue}{01010110} = 10001011$
  \item $10001011 \ll 1 = 00010111$
\end{enumerate}

Jadi, $\text{C1} = 00010111$. Cara yang sama dilakukan untuk $\text{P2}$ dan $\text{P3}$.

Dekripsi blok cipherteks pertama, $00010111$, dengan fungsi $\text{D}$ sebagai berikut:
\begin{enumerate}
  \item $00010111 \gg 1 = 10001011$
  \item $10001011 \oplus \text{K} = 10001011 \oplus \textcolor{blue}{01010110} = 11011101 = \text{P1}$
\end{enumerate}

\end{document}
